\chapter{Nilpotent orbits}
\section{Adjoint and coadjoint orbits}
\subsection{Review on Lie groups and infinitesimal actions}
\begin{itemize}
    \item Let $F:M\rightarrow N$ be a diffeomorphism and $X\in\fr{X}(M)$. We define the pushforward 
    $F_*X\in\fr{X}(N)$ to be vector field defined by the equation $F_*X |_{F(p)}=F_*X_p$. It's a general fact that 
    if $X,Y\in\fr{X}(M)$, then $F_*[X,Y]=[F_*X,F_*Y]$.\par 
    Moreover, if $X\in\fr{X}(M)$ and $Y\in\fr{X}(N)$, we say that $X$ and $Y$ are $F-$related (and we will write 
    $X\sim_FY$) if $F_*X=Y$.
    \item Let $G$ be a Lie group. Define $\fr{X}_{\text{inv}}(G)=\{X\in\fr{X}(G):\ell_{g*}X=X\}$ to be the set of 
    left invariant vector fields of $G$. Then, it's clear that 
    $[\fr{X}_{\text{inv}}(G),\fr{X}_{\text{inv}}(G)]\subset \fr{X}_{\text{inv}}(G)$, so $\fr{X}_{\text{inv}}(G)$ is a 
    Lie algebra. We have an isomorphism 
    \begin{equation}
        \fr{g}\overset{\cong}{\to}\fr{X}_{\text{inv}}(G)=T_eG: X\mapsto [g\mapsto \ell_{g,*e}X]
    \end{equation}
    As a matter of notation, we will write $u^L$ for the left invariant vector field generated by $u\in\fr{g}$, i.e.,
    $u^L(g)=\ell_{g,*}u$.
    \item A $G$- manifold $M$ is a triple $(M,G,\psi)$ where $G$ is a Lie group and 
    $\psi:G\times M\to M$ is a smooth map such that $\psi_g=\psi(g,\cdot)$ defines a homomorphism 
    $G\to\text{Diff}(M):g\mapsto \psi_g$. An infinitesimal generator of $\psi$ is a vector field defined by 
    \begin{equation}
        u_M(x)=\frac{d}{dt}\bigg|_{t=0}\psi_{\exp(tu)}(X).
    \end{equation}
    This defines a map $\fr{g}\to\fr{X}(M):u\mapsto u_M$ which we will show to be a Lie algebra anti-homomorphism.
    \begin{example}
        (infinitesimal generator of $\Ad$). Consider the adjoint action $G\times\fr{g}\to\fr{g}$. We have 
        \begin{equation}
            \begin{split}
                u_{\fr{g}}(v)&=\frac{d}{dt}\bigg|_{t=0}\Ad_{\exp(tu)}(v)\\
                &=\frac{d}{dt}\bigg|_{t=0}(r_{\exp(-tu)}\circ \ell_{\exp(tu)})_{*,e}(v)\\
                &=\frac{d}{dt}\bigg|_{t=0}d_{\exp(tu)}r_{\exp(-tu)}(v^L(\exp(tu)))\\
                &=\frac{d}{dt}\bigg|_{t=0} r_{\exp(tu)*}(v^L)\big|_e\\
                &=\mathcal{L}_{u^L}(v^L)\big|_e\\
                &=[u^L,v^L]_e\\
                &=[u,v].
            \end{split}
        \end{equation}
        Therefore $\Ad_{*,e}=\ad$.
    \end{example}
    \item We have three important results that we will need: 
        \begin{itemize}
            \item[1.] Let $\varphi:M_1\to M_2$ be a $G$ equivariant map, i.e., $\varphi(gx)=g\varphi(x),\;\forall x\in M_1$, then 
            for any $u\in\fr{g}$, we have $u_{M_1}\sim_\varphi u_{M_2}$.
            \item[2.] Let $M$ be a $G$-manifold, then $\fr{g}\to\fr{X}(M):u\mapsto u_M$ is a Lie algebra anti-homomorphism. 
            In fact, the result is true for $(G,\psi(g,\cdot)=r_{g^{-1}})$. This follows frmo the equation 
            \begin{equation}
                u_G(g)=\frac{d}{dt}\bigg|_{t=0}r_{\exp(-tu)}(g)=-u^L(g).
            \end{equation}
            Remember that $r_{\exp(tu)}$ defines the flux of the left invariant vector field $u^L$. Now, we observe that 
            the result is also true for $\overline{M}=G\times M$ with action $g\cdot(a,x)=(r_{g^{-1}}(a),x)$. In fact 
            \begin{equation}
                \begin{split}
                    \pi_{G*(a,x)}u_{\overline{M}}(a,x)&=\frac{d}{dt}\bigg|_{t=0}\pi_G\circ(r_{\exp(-tu)}(a),x)\\
                    &=\frac{d}{dt}\bigg|_{t=0}r_{\exp(-tu)}(a)\\
                    &=-u^L(a).
                \end{split}
            \end{equation}
            Therefore
            \begin{equation}
                [u_{\overline{M}},v_{\overline{M}}]=
                -([u^L,v^L],0)=
                -([u^L,v^L],0)=
                -[u,v]_{\overline{M}}.
            \end{equation}
            Now define the $G$-equivariant map $F:\overline{M}\to M:(a,x)\mapsto a^{-1}x$. By applying $F_*$ to the previous 
            equation we get $[u_M,v_M]=-[u,v]_M$.
            \item[3.] We have $g_*(u_M)=\Ad_g(u)_M$. This is true in $\overline{M}=G\times M$. In fact, we have 
            \begin{equation}
                \begin{split}
                    \pi_G*(a,x)g_{*,(a,x)}u_{\overline{M}}(a,x)
                    &=(\pi_G\circ\psi_g)_{*,(a,x)}u_{\overline{M}}(a,x)\\
                    &=\frac{d}{dt}\bigg|_{t=0}\pi_G\circ\psi_g\circ(r_{\exp(-tu)}(a),x)\\
                    &=\frac{d}{dt}\bigg|_{t=0}\pi_G(r_{g^{-1}}r_{\exp(-tu)}(a),x)\\
                    &=\frac{d}{dt}\bigg|_{t=0}a\exp(-tu)g^{-1}\\
                    &=\frac{d}{dt}\bigg|_{t=0}ag^{-1}c_g(\exp(-tu))\\
                    &=\frac{d}{dt}\bigg|_{t=0}ag^{-1}\exp(-\Ad_g(u))\\
                    &=\frac{d}{dt}\bigg|_{t=0}r_{\exp(-\Ad_g(u))}(ag^{-1})\\
                    &=-(\Ad_g(u))^L(ag^{-1}).
                \end{split}
            \end{equation}
            Finally 
            \begin{equation}
                \begin{split}
                    g_*u_{\overline{M}}\big|_{g\cdot(a,x)}&=-(\Ad_g(u)^L(ag^{-1}),0)\\
                    &=-(\Ad_g(u)^L,0)\big|_{g\cdot(a,x)}.
                \end{split}
            \end{equation}
        \end{itemize}
        Now we have 
        \begin{equation}
            g_*u_M=g_*F_*u_{\overline{M}}=F_*g_*u_{\overline{M}}=F_*(\Ad_g(u))_{\overline{M}}=(\Ad_g(u))_M.
        \end{equation}
        \item Let $(M,G,\psi)$ be a $G$-manifold. For $x\in M$, we define the isotropy subgroup 
        $G_x=\{g\in G:gx=x\}$. It's clear that $G_x\subset G$ is closed, so it's a Lie subgroup of $G$. Define 
        $\fr{g}_x=\{u\in\fr{g}:u_M(x)=0\}$, we claim that  
        \begin{equation}
            \text{Lie}(G_x)=\fr{g}_x.
        \end{equation}
        Let $\exp(tu)\in G_x,\;\forall t\in\bbR$, that is, assume that $u\in\text{Lie}(G_x)$. Then, by definition, we have 
        $\exp(tu)x=x$. So, by differentiating this equation we get $u_M(x)=0$ and $u\in\fr{g}_x$. Conversely, assume that 
        $u\in\fr{g}_x$. Let $g=\exp(ts),\;s\in\bbR$, then 
        \begin{equation}
            \frac{d}{dt}\bigg|_{t=s}\exp(tu)x= g_{*,x}(\frac{d}{dt}\bigg|_{t=0}\exp(tu)x)=0
        \end{equation}
        and $\exp(tu)\cdot x = x,\forall t\in\bbR$. Therefore $u\in\text{Lie}(G_x)$.
        \item Now, if $\mathcal{O}_x=Gx$, it's a fact that if we identify $G/G_x\cong \mathcal{O}_x$ by the 
        orbit stabilizer theoerem, then $\mathcal{O}_x$ inherits a differentiable structure such that the inclusion 
        $i:\mathcal{O}_x\hookrightarrow M$ is an imersion. In this context we claim that 
        \begin{equation}
            i_{*,x}T_x\mathcal{O}_x=\{u_M(x): u\in\fr{g}\}.
        \end{equation}
        Lets prove it. We know that $i:G/G_x\hookrightarrow M:[g]\mapsto gx$ is an imersion. First observe that 
        \begin{equation}
            \begin{split}
                \frac{d}{dt}\bigg|_{t=0}i\circ\pi\circ r_{\exp(tu)}(g)&=\frac{d}{dt}\bigg|_{t=0}g\exp(tu)\cdot x\\
                &=g_{*,x}u_M(x)\\
                &=(\Ad_g(u))_M(gx).
            \end{split}
        \end{equation}
        Since $\pi:G\rightarrow G/G_x$ is a submersion, we conclude that 
        $i_{*,[g]}T_{[g]}(G/G_x)\subset \{u_M(g\cdot x):u\in\fr{g}\}$. Now define the map 
        \begin{equation}
            \fr{g}\to \{u_M(g\cdot x):u\in\fr{g}\}
        \end{equation}
        which has kernel equal to $\text{Lie}(G_{g\cdot x})$. Therefore we have an isomorphism
        \begin{equation}
            \fr{g}/\text{Lie}(G_{g\cdot x})\overset{\cong}{\to}\{u_M(g\cdot x):u\in\fr{g}\}.
        \end{equation}
        Therefore the dimension of the last subspace is equal to 
        \begin{equation}
            \begin{split}
                \dim(G)-\dim(G_{g\cdot x})&=\dim(G)-\dim(gG_xg^{-1})\\
                &=\dim(G)-\dim(G_x)\\
                &=\dim(G/G_x).
            \end{split}
        \end{equation}
        Therefore, by dimensional reasons we have $i_{*,[g]}T_{[g]}(G/G_x)=\{u_M(g\cdot x):u\in\fr{g}\}$ and since 
        $i_*$ is injective we can identify $T_x{\mathcal{O}_x}=\{u_M(x):u\in\fr{g}\}$.
\end{itemize}

\subsection{Co-adjoint orbits and KKS form}

    Let $G$ be a Lie group. We will show that the orbits on $\fr{g}^*$ under the co-adjoint action have a natural
    structure of a symplectic manifold.\par
    \textbf{KKS}: Let $\xi\in\mathcal{O}$, we define $\omega\in\bigwedge^2T_\xi^*\mathcal{O}$ by the formula: 
    \begin{equation}
        \omega_\xi(u_{\fr{g}^*}(\xi),v_{\fr{g}^*}(\xi))=\langle\xi,[u,v] \rangle.
    \end{equation}
    We need to verify a few things:
\begin{itemize}
    \item $\omega_\xi$ is well defined. Assume that $u_{\fr{g}^*}(\xi)=u_{\fr{g}^*}^\prime(\xi)$ and 
    $v_{\fr{g}^*}(\xi)=v_{\fr{g}^*}^\prime(\xi)$, then 
    \begin{equation}
        \begin{split}
            \omega_\xi(u_{\fr{g}^*}^\prime(\xi),v_{\fr{g}^*}^\prime(\xi))&=
            -u_{\fr{g}^*}^\prime(\xi)(v^\prime)\\
            &=-u_{\fr{g}^*}(\xi)(\xi)(v^\prime)\\
            &=\langle \xi,[u,v^\prime] \rangle\\
            &=v_{\fr{g}^*}^\prime(\xi)(u)\\
            &=v_{\fr{g}^*}(\xi)(u)\\
            &=\langle\xi,[u,v] \rangle.
        \end{split}
    \end{equation}
    \item It's skew symmetric and nondegenerate (clear): $\langle\xi,[u,v] \rangle=0\;\forall v\iff u_{\fr{g}^*}(\xi)=0$.
    \item It's $G-$invariant, so it's smooth. This can be verified directly by using the relation $g_*(u_{\fr{g}^*})=(\Ad_g(u))_{\fr{g}^*}$.
    \item By the previous item we see that $\psi_{e^{tu}}^*\omega=\omega$. Therefore 
    \begin{equation}
        \mathcal{L}_{\fr{g}^*}\omega=0\iff (d\iota_{u_{\fr{g}^*}}+\iota_{u_{\fr{g}^*}}d)\omega=0.
    \end{equation}
    \textbf{Claim}: We have $\iota_{u_{\fr{g}^*}}\omega=du$ where we identify $u\in\fr{g}\cong C^\infty_{\text{Lin}}(\fr{g}^*)$.\par 
    This can be seen directly by the equation:
    \begin{equation}
        \begin{split}
            \iota_{u_{\fr{g}^*}}\omega\big|_\xi(v_{\fr{g}^*}(\xi))&=\langle\xi,[u,v] \rangle\\
            &=\langle\xi,[u,v] \rangle\\
            &=v_{\fr{g}^*}(\xi)(u)\\
            &=du\big|_\xi v_{\fr{g}^*}(\xi)\;(\text{by the identification }\fr{g}^{**}=\fr{g})
        \end{split}
    \end{equation}
    \item For $G=SO(3)$ the co-adjoint and adjoint orbits are the same. We have an isomorphism of Lie algebras
    $\bbR^3\cong\fr{so}(3)$ which maps $e_i\in\bbR^3$ to the associated matrix representation of the map 
    $(e_i)_\times:v\mapsto e_i\times v$. Moreover, we have the identity 
    \begin{equation}
        A(u\times v)= Au\times Av,\; \forall A\in SO(3).
    \end{equation}
    With this equation we get $(Au)_\times v = Au\times v = Au_\times A^{-1}(v)=\Ad_A(u_\times)v$. So this isomorphism 
    is $SO(3)-$equivariant, with the adjoint orbits being spheres in $\bbR^3$.
    \begin{dang}
     \begin{remark}
        If $\nu: \fr{g}\to \fr{g}^*$ is an isomorphism that intetwine the adjoint and co-adjoint actions of $G$. Then we 
        can induce a symplectic form on adjoint orbits using KKS. The formula is the one 
        \begin{equation}
            \omega_\xi(u_{\fr{g}}(\xi),v_{\fr{g}}(\xi))=\langle \nu(\xi), [u,v]\rangle.
        \end{equation}
    \end{remark}    
    \end{dang}

    Going back to definitions: If $\xi\in(\bbR^3)^*$, we have 
    \begin{equation}
        \begin{split}
            \omega_\xi(u_{\fr{g}}(\xi),v_{\fr{g}}(\xi))&=\langle\xi,[u,v]\rangle\\
        &=\langle \xi, u\times v\rangle \;(\bbR^3\text{ is a Lie algebra with the croos product})\\
        &=\xi\cdot (u\times v)\;(\text{By the isomorphism }\bbR^3\cong(\bbR^3)^*)
        \end{split}
    \end{equation}
    \begin{dang}
    \begin{remark}
        Let $S\subset M$ a hypersurface, $M$ oriented with orientation form $\omega$. Let 
        $N\in\fr{X}(M)$ be a vector field nowhere tangent to $S$, then $\iota_N\omega|_S$ is an orientation 
        form for $S$.
    \end{remark}        
    \end{dang}

    By the previous remark we can use the vector field 
    $N= x\tfrac{\partial}{\partial x}+y\tfrac{\partial}{\partial y}+z\tfrac{\partial}{\partial z}$ to induce 
    an orientation form 
    \begin{equation}
        \sigma = xdy\wedge dz - ydx\wedge dz+zdx\wedge dy \big|_{\mathbb{S}^2}
    \end{equation}
    on the spheres. In fact, if $f:(x,y,z)\mapsto x^2+y^2+z^2-1$, then $df(N)=2(x^2+y^2+z^2)$ is never zero on $\mathbb{S}^2$, so 
    $N$ is nowhere tangent to $\mathbb{S}^2$.\par 

    \begin{dang}
        \begin{remark}
            If $f:\bbR^n\to\bbR$ is smooth with regular value $c$, then $S=f^{-1}(0)$ is a regular submanifold 
            with tangent spaces 
            \begin{equation}
                i_{*p}T_pS \cong \nabla f(p)^\perp
            \end{equation}
        \end{remark}
    \end{dang}

    In our case we have $\nabla f(\xi) = \xi$, so $T_\xi\mathbb{S}^2=\xi^\perp=\{\xi\times u:u\in\bbR^3\}$. If we 
    use the area form 
    \begin{equation}
        \begin{split}
            \sigma_\xi(\xi\times u,\xi\times v) &=\xi\cdot ((\xi\times u)\times (\xi\times v))\;(\text{use the area form})\\
            &=-(\xi\times u)\cdot (\xi\times(\xi\times v))\\
            &=-(\xi\times u)\cdot (\xi(\xi\cdot v)-v(\xi\cdot\xi))\\
            &=\|\xi\| (\xi\times u)\cdot v \\
            &=\xi\cdot (u\times v)\|\xi\|.
        \end{split}
    \end{equation}

    In conclusion, we see the relation between the KKS and the area form:
    \begin{equation}
        \omega_{\text{kks}}=\frac{1}{r}\sigma
    \end{equation}
    where $r=\|\xi\|$ is the radius of the sphere.
\end{itemize}

\subsection{Hamiltonian actions}
Let $X\in T_0\bbR^n\cong \bbR^n$, then $\widetilde{X}(g)=(g,X)$. In fact
\begin{equation}
    \begin{split}
         \ell_{g*,0}(\frac{\partial}{\partial x_i}\bigg|_g)x_j &= \frac{\partial}{\partial x_i}\bigg|_g x_j\circ \ell_g\\
         &=\frac{\partial}{\partial x_i}\bigg|_g(x_j+x_j(g))=\delta_{ij}.
    \end{split}
\end{equation}

In general if $\widetilde{X}\in\fr{X}_{\text{inv}}(G)$, then we set $\exp(X)=\gamma_1^{\widetilde{X}}(e)$. If $G=\bbR^n$, then 
$\gamma_t^{\widetilde{X}}(0)=tX$ and $e^X=X$.


    Let $(M,\omega)$ be a symplectic manifold. Then

    \begin{itemize}
        \item $G-$action is symplectic if $\mathcal{L}_X\omega = 0$. This condition is equivalent to 
        $\psi_g^*\omega=\omega\;\forall g\in G$ if $G$ is connected.
        \item $G-$action is Hamiltonian if there exist a linear map $\hat{\mu}:\fr{g}\to C^\infty(M)$
        such that the following diagram commutes: 
        \begin{center}
            \begin{tikzcd}
                                             &   & C^\infty(M)\arrow[dd, dashed]& \hat{\mu}(u)\arrow[dd, mapsto] \\
    \fr{g}\arrow[rru, "\hat{\mu}"]\arrow[rrd]&u\arrow[rru, mapsto]\arrow[rrd, mapsto]  &    &\\
                                             &   & \fr{X}(M)   &u_M = X_{\hat{\mu(u)}}
            \end{tikzcd}
        \end{center}
    \end{itemize}

    \begin{remark}
        We have a bijection 
        \begin{equation}
            \{\hat{\mu}: \fr{g}\to C^\infty(M)\} \leftrightarrow \{\mu:M\to\fr{g}^*\}
        \end{equation}
    \end{remark}

\begin{itemize}
    \item The action is weakly Hamiltonian if
    \begin{equation}
        u_M = X_{\hat{\mu}(u)} \iff \iota_{u_M}\omega = d\hat{\mu} = d\langle\mu,u\rangle.
    \end{equation}
    \item The action is Hamiltonian if, in addition, we have 
    \begin{equation}
        \mu \circ \psi_g = \Ad_g^*\circ\mu \overset{G\text{ connected}}{\iff} \hat{\mu}([u,v])=\{\hat{\mu}(u),\hat{\mu}(v)\}.
    \end{equation}
    The relation between the two definitions is $\hat{\mu}(u)\big|_x = \langle\mu(x),u\rangle$.
\end{itemize}

Now, look at the following equation for $x\in M, v\in \fr{g}$:
\begin{equation}
    \langle \mu\circ\psi_g(x),v\rangle = \langle\Ad_g^*\circ\mu(x),v\rangle = \langle \mu(x), \Ad_{g^{-1}}(v)\rangle.
\end{equation}
This equation means that $\hat{\mu}(v)\circ\psi_g(x)=\hat{\mu}(\Ad_{g^{-1}}(v))\big|_x$. If we put $g=e^{tu},u\in\fr{g}$ and 
differentiate this equation, we get $\hat{\mu}([u,v])=\{\hat{\mu}(u),\hat{\mu}(v)\}$. For the converse we need the hypothesis 
of $G$ being connected.\par 

\begin{remark}
    Remember that $\fr{X}_{\text{Ham}}(M)\subset\fr{X}_{\text{Sym}}(M)$. 
\end{remark}

In our case we have $\mathcal{L}_{u_M}\omega = 0$. Since $u_M$ has flux given by $\psi_{e^{tu}}$, we get 
$\psi_{e^{tu}}^*\omega=\omega$. Therefore $\psi_g^*\omega = \omega,\;\forall g\in G^\circ$.
    
\begin{dang}
    \begin{theorem}
        (Noether principle). Let $(M,\omega,\mu)$ be a Hamiltonian $G$-space and $H\in C^\infty(M)$. Then 
        $H$ is $G$-invariant if and only if $\mu$ is preserved by the Hamiltonian flux of $H$.
    \end{theorem}
\end{dang}
Proof: Remember that $H$ is $G$ invariant if $\psi_g^*H=H\iff\mathcal{L}_{u_M}H=0$ if $G$ is connected. Also, 
$\mu$ is preserved by the Hamiltonian flux of $H$ if $\varphi_t^*\mu=\mu$ where $\varphi_t=\varphi_t^{X_H}$ is the 
flux of $X_H$. This equation is true if 
\begin{equation}
    \begin{split}
        \mathcal{L}_{X_H}\langle\mu,u\rangle &= \frac{d}{dt}\bigg|_{t=0}\varphi_t^*\langle\mu,u\rangle\\
           &=\frac{d}{dt}\bigg|_{t=0}\langle\varphi_t^*\mu,u\rangle=0\;\forall u\in\fr{g}\\
           &\iff \varphi_t^*\mu =\mu,\;\forall t\in\bbR.
    \end{split}
\end{equation}
But we see that 
\begin{equation}
    \begin{split}
        \mathcal{L}_{X_H}\langle\mu,u\rangle &= (d\iota_{X_H}+\iota_{X_H}d)\langle\mu,u\rangle\\
        &=\iota_{X_H}d\langle\mu,u\rangle\\
        &=\iota_{X_H}\iota_{u_M}\omega\\
        &=-dH(u_M)\\
        &=-\mathcal{L}_{u_M}H.
    \end{split}
\end{equation}
\qed

\begin{example}
    A general situation: $(M,\omega)$ is symplectic with $\omega = -d\theta$ exact. Suppose that we have 
    a $G$ on $M$ such that $\psi_g^*\omega = \omega$. In this case we define 
    \begin{equation}
        \iota_{u_M}\theta =  \langle \mu,u\rangle.
    \end{equation}
    It's easy to see that this defines a Hamiltonian action in $M$.
\end{example}

\begin{example}
    (Cotangent lifting). Take a $G$-action $\psi_g$ on $M$, then we can lift 
    it to a $G$-action $\widehat{\psi}_g$ on $T^*M$. So we have $\widehat{\psi}_g^*\alpha=\alpha$ and thus we get 
    a Hamiltonian map given by $\langle\mu,u\rangle = \iota_{u_M}\alpha$. 
\end{example}

\begin{example}
    Let $G=\bbR^n$ acting on $M=\bbR^n$ by translation $g\cdot q = q+g$. With the identification 
    $T^*M\cong\bbR^n\times\bbR^n$ we have $g\cdot(q,p)=(g+q,p)$. We can explictly compute the moment map:
    \begin{equation}
        \begin{aligned}
            &u = (a,b) \mapsto u_M = \sum_ia_i\frac{\partial}{\partial q_i}+\sum_ib_i\frac{\partial}{\partial p_i}\\
            &\iota_{u_M}\sum_ip_idq_i=\sum_ip_idq_i(u_M)=\sum_ip_ia_i\\
            &\mu(q,p):\bbR^n\to\bbR^n:e_i\mapsto p_i 
        \end{aligned}
    \end{equation}
    So $\mu(q,p)=p$ in the choosen basis.
\end{example}

\begin{example}
    Let $G=SO(3)$ acting in $\bbR^3$ in the usual way a lift it to an action on $T^*\bbR^3=\bbR^3\times\bbR^3$. 
    Now observe that 
    \begin{equation}
    \begin{split}
    \psi_{g*,v}(\frac{\partial}{\partial x_i}\bigg|_v)x_j&=\frac{\partial}{\partial x_i}\bigg|_vx_j\circ \psi_g \\
    &=d(x_j\circ \psi_g)_v(e_i)\\
    &=x_j \circ\psi_g(e_i)\\
    &=g_{ji}
    \end{split}
    \end{equation}
    This proves that $\psi_{g*,v}(\sum_ia_i\frac{\partial}{\partial x_i}\bigg|_v)
    =\sum_i(ga)_i\frac{\partial}{\partial x_i}\bigg|_{gv}$. Therefore, the corresponding action on $T\bbR^3$ is 
    \begin{equation}
        g\cdot(q,p)=(gq,gp).
    \end{equation}
    Remember that the inverse of the dualization preserve base points and inverse the dual map, so the action on 
    $T^*\bbR^3$ is 
    \begin{equation}
        g\cdot (q,p)=(gq,(g^{-1})^Tp)=(gq,gp),\;g\in SO(3).
    \end{equation}
    Finnaly we see that 
    \begin{equation}
        \begin{aligned}
            &u_M(q,p)=\frac{d}{dt}\bigg|_{t=0}(e^{tu}q,e^{tu}p)=(uq,up)\\
            &u_M = \sum_i (uq)_i\frac{\partial}{\partial q_i}+\sum_i (up)_i\frac{\partial}{\partial p_i}\\
            &\iota_{u_M}\sum_ip_idq_i = \sum_i p_iu_{ij}q_j\\
            &\mu(q,p):\fr{so}(3)\to \bbR: \ell_1\mapsto p_3q_2-p_2q_3,\dots\Rightarrow \mu(q,p)=q\times p
        \end{aligned}
    \end{equation}
\end{example}
\subsection{Marsden, Weinstein, Meyer theorem}

\subsection{Lazy and sophisticated approach for coadjoint orbits}
\begin{dang}
    \begin{remark}
        Let $\psi$ be a nondegenerate, symmetric bilinear form on a vector space $E$. If $W\subset E$ is a subspace 
        then $E=W\oplus W^\perp$.\par
        In fact, if $\nu:E\to E^*$ is the linear map induced by $\psi$, then we see that $\nu(W)=\text{Ann}(W)$. Moreover,
        the exact sequence 
        \begin{equation}
            0\to \text{Ann}(W)\to E^*\to W^*\to0
        \end{equation}
        yields $\dim(\text{Ann}(W))+\dim(W)=\dim(E)$. The dimension formula yields $\dim(W)+\dim(W^\perp)=\dim(E)$, so 
        $W\cap W^\perp=(0)$. Putting together all these results, we get an isomorphism 
        $W\oplus W^\perp\to E:(w,w^\prime)\mapsto w+w^\prime$.
    \end{remark}
\end{dang}

\begin{dang}
    \begin{lemma}
        Let $\fr{g}=\fr{c}\oplus [\fr{g},\fr{g}]$ be a reductive Lie algebra and 
        $\psi$ a symmetric, nondegenerate, invariant bilinear form on $\fr{g}$. 
        Then, for any two simple components $\fr{s}_1,\fr{s}_2\subset[\fr{g},\fr{g}]$ we have 
        $\psi(\fr{s}_1,\fr{s}_2)=0$, $\psi(\fr{s}_1,\fr{c})=0$.
    \end{lemma}
\end{dang}
Proof: We just observe that $[\fr{s}_1,\fr{s}_2]=0$, while $[\fr{s},\fr{s}]=\fr{s}$ for $\fr{s}$ simple. Therefore 
\begin{equation}
    \psi(\fr{s}_1,\fr{s}_2)=\psi([\fr{s}_1,\fr{s}_1],\fr{s}_2)=\psi(\fr{s}_1,[\fr{s}_1,\fr{s}_2])=0.
\end{equation}
In the same way we see that $\psi(\fr{s},\fr{c})=0$.\qed

\begin{dang}
    \begin{lemma}\label{Lemma: lazy vs sophisticated approach to nilpotent orbits}
        (1.3.9 - Lazy and sophisticated approach):
        Let $\fr{g}$ be a reductive Lie algebra with nondegenerate, symmetric, invariant bilinear form $\psi$. 
        If $f=\psi(X,\cdot)\in\fr{g}^*$, then 
        \begin{equation}
            f|_{\fr{g}^f}=0\iff X\text{ is nilpotent}.
        \end{equation}
    \end{lemma}
\end{dang}
Proof: First assume $f=\psi(X,\cdot)$ with $X$ nilpotent, i.e., $X\in[\fr{g},\fr{g}]$ and $\ad(X)$ is a nilpotent
endomorphism of $\fr{g}$. Now, let $[\fr{g},\fr{g}]=\fr{s}_1\oplus\cdots\oplus\fr{s}_r$ and 
$X=X_1+\cdots+X_r, Z=Z_0+Z_1+\cdots+Z_r$ for $Z_0\in\fr{c},\;X_i,Z_i\in\fr{s}_i$. 


Then 
\begin{equation}
    \begin{split}
        f(Z)=\psi(X,Z)&=\sum_{i,j}\psi(X_i,Z_j)\\
        &=\sum_i\psi(X_i,Z_i)\\
        &=\sum_i\alpha_i K^{\fr{s}_i}(X_i,Z_i)\\
        &=\sum_{i,j}\alpha_j K(X_i,Z_j)\\
        &=K(X,Z^\prime),\;Z^\prime = \sum_j\alpha_j Z_j\\
        &=\Tr(\ad(X)\circ\ad(Z^\prime)),\;Z^\prime = \sum_j\alpha_j Z_j.
    \end{split}
\end{equation}
However, if $Z\in\fr{g}^f=\fr{g}^X$, then $[Z,X]=0\Rightarrow [Z_i,X_i]=0\;\forall i$ and $Z^\prime \in\fr{g}^X$ as well. 
In fact, we have $Z\in\fr{g}^X\iff Z^\prime \in\fr{g}^X$ because $\alpha_i\neq0,\;\forall i$.
Then $[\ad(X),\ad(Z^\prime)]=0$ and since $\ad(X)$ is nilpotent, we see that $\ad(X)\circ\ad(Z^\prime)$ is nilpotent and thus
must be of trace zero.\par


Now assume that $f|_{\fr{g}^f}=0$. We need to show two things: That $X\in[\fr{g},\fr{g}]$ and $\ad(X)$ is nilpotent. 
For the first one we have $f|_{\fr{g}^f}=\psi(X,\cdot)|_{\fr{g}^X}=0$, so $X\in(\fr{g}^X)^\perp$. Moreover, it's clear that 
$\psi([X,\fr{g}],\fr{g}^X)=0$ by invariance. So $[X,\fr{g}]\subset (\fr{g}^X)^\perp$, and since the two subspaces have the same dimension, we get 
$[X,\fr{g}]=(\fr{g}^X)^\perp$. So we see that $X\in(\fr{g}^X)^\perp=[X,\fr{g}]$.\par 

Now we need to show that $\ad(X)$ is nilpotent. To see this, we use the fact that $X=[H,X]$ for some $H\in\fr{g}$. 
Now, decompose $\fr{g}$ into generalized eigenspaces:
\begin{equation}
    \fr{g}=\bigoplus_{s\in\bbC}E_s^{\ad(H)}=\bigoplus_{s\in\bbC}\fr{g}_s.
\end{equation}
We claim that if $W\in \fr{g}_s$, then $\ad(X)W\in \fr{g}_{s+1}$. To see this, apply 

\begin{equation}
    \begin{split}
        (\ad(H)-s-1)^N\ad(X)W&=\sum_{k=0}^s\binom{N}{k}\ad(H)^k\ad(X)W(-s-1)^{N-k}\\
        &=\sum_{k=0}^N\binom{N}{k}\ad(X)(\ad(H)+1)^kW(-s-1)^{N-k}\\
        &=\ad(X)(\ad(H)+1-s-1)^NW\\
        &=\ad(X)(\ad(H)+1)^NW=0.
    \end{split}
\end{equation}
\qed

\begin{dang}
    \begin{remark}
        Since $\psi$ is nondegenerate, we have the decomposition at the level of vector spaces:
        \begin{equation}
            \fr{g}=\fr{g}^X\oplus(\fr{g}^X)^\perp.
        \end{equation}
        Now, by dimension formula we see that $\dim(\fr{g}^X)^\perp=\dim(\fr{g})-\dim(\fr{g}^X)=\dim([X,\fr{g}])$.
        Since $K([X,\fr{g}],\fr{g}^X)=0$ by invariance, see that $[X,\fr{g}]\subset(\fr{g}^X)^\perp$ and the equality
        follows by dimensional reasons.
    \end{remark}
\end{dang}
\subsection{Standard triples}
For each $r\geq1$ we have an irreducible representation of $\fr{sl}_2$ on $\bbC^{r+1}$ given by:

  \begin{equation}
    \begin{split}
        &\rho_r(X)=\begin{pmatrix}
0&1&&&\\ &0&\rotatebox{-45}{$\cdots$}&&\\ &&\rotatebox{-45}{$\cdots$}&\rotatebox{-45}{$\cdots$}&\\&&&0&1\\&&&&0
\end{pmatrix},\\
        &\rho_r(Y)=\begin{pmatrix}
    0&&&&\\ \mu_1&0&&&\\ &\rotatebox{-45}{$\cdots$}&\rotatebox{-45}{$\cdots$}&&\\ &&\rotatebox{-45}{$\cdots$}&0&\\ &&&\mu_r&0
\end{pmatrix},\;\mu_i=i(r+1-i)\\
        &\rho_r(H)=\text{Diag}(r,r-2,\dots,-r+2,-r).
    \end{split}
  \end{equation}
In particular this defines a homomorphism $\fr{sl}_2\to\fr{sl}_{r+1}$. So, to each nilpotent orbit 
    $\mathcal{O}=\mathcal{O}_{d_1,\dots, d_r}$ in $\fr{sl}_n\;(n=d_1+\cdots+d_r)$, we can find a nilpositive element 
$X\in\mathcal{O}$. In fact, if $k=l([d_1,\dots,d_r])$ we put 
$\rho(X)=\bigoplus_{i=1}^k\rho_{d_i-1}(X)=X_{d_1,\dots, d_r}\in\fr{sl}_n$ and 
$\{\rho(X),\rho(H),\rho(Y)\}$ is a $\fr{sl}_2-$triple.



\section{Jacobson-Morozov theorem}

\begin{dang}
    \begin{theorem}
        Let $X\in\fr{g}$ be a nilpotent element of a complex semisimple Lie algebra. Then $X$ is the nilpositive element
        of some $\fr{sl}_2-$triple in $\fr{g}$.
    \end{theorem}
\end{dang}
Proof: The proof will be divided in four steps:
\begin{itemize}
    \item Induction on dimension: The base case if $\dim(\fr{g})=3$, i.e., $\fr{g}=\fr{sl}_2$. If $X\in\fr{g}$ is nilpotent, then 
    we see that there exists $g\in GL_n$ such that 
    $\Ad(g)X=\left(\begin{smallmatrix}0&1\\0&0\end{smallmatrix}\right)$. Putting $g\to g/\sqrt{\det(g)}$, we see that we have 
    in fact an element $g\in G=SL_2$ such that $\Ad(g)X=\left(\begin{smallmatrix}0&1\\0&0\end{smallmatrix}\right)$. So 
    we have an automorphism $\phi\in\Aut(\fr{g})$ such that $\phi(X)$ is the nilpositive element of the standard triple 
    in $\fr{sl}_2$. Now, we will proceed the proof by assuming $\dim(\fr{g})>3$ and that 
    $X$ doesn't lie in any proper semisimple Lie algebra of $\fr{g}$.
    \item By the proof of theorem~(\ref{Lemma: lazy vs sophisticated approach to nilpotent orbits}), we see that 
    $K(X,\fr{g}^X)=0$ (here, since $\fr{g}$ is semisimple, we can just take $\psi=K$). Therefore, $X\in(\fr{g}^X)^\perp$
    which is equal to $[X,\fr{g}]$ and there exists $H\in\fr{g}$ such that $[H,X]=2X$. \par
    In this condition, we can assume $H$ to be semisimple: Since $\ad(H)$ restricts to $\bbC X$ and $\ad(H)_{s,n}$ are 
    polynomials in $\ad(H)$, then its semisimple and nilpotent parts also restricts to $\bbC X$. In particular 
    $\ad(H_n)|_{\bbC X}^{\dim (\bbC X)}=0$, i.e., $\ad(H_n)(X)=0$. Therefore $[H,X]=[H_s,X]$ and we can pass notation to 
    use $H$ as $H_s$.
    \item In this part of the proof we will use the induction step to show that $K(H,\fr{g}^X)=0$, i.e., 
    $H\in(\fr{g}^X)^\perp=[X,\fr{g}]$.\par 

    Suppose that $H\notin (\fr{g}^X)^\perp$, so $K(H,\fr{g}^X)\neq0$. We know that $\fr{g}^X$ is $\ad(H)-$invariant:
    \begin{equation}
        \ad(X)\ad(H)Y=(\ad(H)\ad(X)+\ad([X,H]))Y=-2\ad(X)Y = 0.
    \end{equation}
    So, since $\ad(H)$ is semisimple, its restriction to $\fr{g}^X$ is also semisimple and we have the eigenspace decomposition
    \begin{equation}
        \fr{g}^X=\bigoplus_{i=1}^k E_{\lambda_i}^{\ad(H)|_{\fr{g}^X}}=\bigoplus_{i=1}^k \fr{g}_{\lambda_i}^X.
    \end{equation}
    In this case we see that if $\lambda_i\neq0$, then 
    \begin{equation}
        K(H,\fr{g}_{\lambda_i}^X)=\frac{1}{\lambda_i}K(H,[H,\fr{g}_{\lambda_i}^X])=0.
    \end{equation}
    So there exist $Z\in\fr{g}_0^X=(\fr{g}^X)^H$ such that $K(H,Z)\neq0$.\par

    If $Z$ is nilpotent, then since $[Z,H]=0$, we see that $K(Z,H)=\Tr\ad(Z)\circ\ad(H)=0$ because 
    $\ad(Z)\circ\ad(H)$ is nilpotent. Therefore, we will assume that $Z_s\neq0$. Here we observe that 
    $Z_s\in(\fr{g}_X)^H$. In fact, whenever $\ad([W,Z])=0$, we have $\ad([W,Z_s])=0$ and $[W,Z_s]=0$.
    \par 
    Finally, we see that since $Z_s\neq0$, then $\fr{g}^{Z_s}\overset{\neq}{\subset}\fr{g}$ and 
    $[\fr{g}^{Z_s},\fr{g}^{Z_s}]$ is a proper semisimple subalgebra of $\fr{g}$ (we are using the fact that 
    centralizers of single elements $\fr{g}^W$ are reductive algebras). However, since $Z_s\in(\fr{g}_X)^H$, we 
    have $H,X \in \fr{g}^{Z_s}$. But $2X=[H,X]\in [\fr{g}^{Z_s},\fr{g}^{Z_s}]$, which contradicts the Induction
    hypothesis.

    \item Now we know that there exist $H\in\fr{g}$ semisimple such that $[H,X]=2X$. Moreover, we know that 
    $H=[X,Y]$ for some $Y\in\fr{g}$. Now, decompose $\fr{g}$ into $\ad(H)-$eigenspaces:
    \begin{equation}
        \fr{g}=\bigoplus_{i=1}^k E_{\lambda_i}^{\ad(H)}=\bigoplus_{i=1}^k\fr{g}_{\lambda_i}.
    \end{equation}
    We see that $\ad(X)\fr{g}_{\lambda_i}\subset \fr{g}_{\lambda_i+2}$. Moreover
    \begin{equation}
        H=[X,Y]=\sum_{i=1}^k[X,Y_i],\; Y_i\in\fr{g}_{\lambda_i}.
    \end{equation}
    Therefore, since $H\in\fr{g}_0$,
     we see that $H=[X,Y_{-2}]$, where $Y_{-2}$ is the eigenvector relative to the eigenvalue $-2$ 
     (we are abusing notation here !). Changing 
     $Y\to Y_{-2}$ we get the standard triple $\{Y,H,X\}$ with $X$ being the nilpositive element fo the standard triple.
\end{itemize}
\qed

\begin{dang}
    \begin{remark}
        In the previous proof we saw that we can take the neutral element $H$ as a semisimple element.
    \end{remark}
\end{dang}

\section{Conjugation theorems}
Some important results to remember:
\begin{itemize}
    \item Let $\{Y,H,X\},\{Y^\prime,H,X\}$ two standar triples with $H=H_s$ being the neutral element and $X=X_n$ the 
    nilpositive element. Then $Y=Y^\prime$. 
    This follows at once by the condition $[Y-Y^\prime,X]=H-H=0$, so $Y-Y^\prime\in\fr{g}^X\cap\fr{g}_{-2}=0$.\qed
    \item Let $X=X_n\neq0$, then $\fr{g}^X\cap[X,\fr{g}]=\bigoplus_{i>0}\fr{g}_i^X$. In fact, if 
    $W\in\fr{g}_i^X\setminus0,\;i>0$, then $W$ is a highest weight for the action of $\fr{sl}_2\cong\bbC\langle Y,H,X\rangle$ 
    (remember the Jacobson Morozov theorem). Since $W$ has weight greater than zero, we necessarily have $Z=[Y,W]\neq0$ because 
    of the representation theory of $\fr{sl}_2$. Therefore $W=\tfrac{1}{i}[X,Z]\in[X,\fr{g}]$. 
    Conversely, if $W=\sum_{i\geq0}W_i,\;W_i\in\fr{g}_i^X$ is also in $[X,\fr{g}]$,
    then $W_0\in[X,\fr{g}]$, because $W_i\in[X,\fr{g}]$ for $i>0$, as we saw before. However, if $W_0=[X,Z]$, then 
    we can assume $Z\in\fr{g}_{-2}$. So if $Z\neq0$, we have $\ad(X)^2Z=[X,W_0]\neq0$ by the representation theory of 
    $\fr{sl}_2$, which is a contradiciton. Then $W_0=0$ and $W\in \bigoplus_{i>0}\fr{g}_i^X$.\qed 
    \begin{dang}
        \begin{notation}
            We will write $\fr{u}^X=\fr{g}^X\cap[X,\fr{g}]=\bigoplus_{i>0}\fr{g}_i^X$. It's clearly a nilpotent ideal in 
            $\fr{g}^X=\bigoplus_{i\geq0}\fr{g}_i^X$.
        \end{notation}
    \end{dang}

    \item Fact: We can integrate $\fr{u}^X$ to a simply connected subgroup $U^X\subset G_{\ad}^X$ (Lie third theorem). 
    Moreover, since $\fr{u}^X$ is nilpotent, the exponential map $\exp:\fr{u}^X\to U^X$ is a \textit{diffeomorphism}. 
    Recall the following calculation:
    \begin{equation}
        \begin{split}
             \Ad_{\ad(Z)}H&=\exp(\ad(Z))H\\
              &=\sum_{i\geq0}\frac{1}{i!}\ad(Z)^iH
        \end{split}
    \end{equation} 
    which is a finite sum, because $\fr{u}^X$ is nilpotent. Moreover $\fr{u}^X=\bigoplus_{i>0}\fr{g}_i^X$ is 
    $\ad(H)-$invariant and $\ad(Z)^iH\in\fr{g}_i^X,\;i>0$. Then $\Ad_{\exp(Z)}H\in H+\fr{u}^X$. We will write 
    \begin{equation}
        \Ad(U^X)H\subset H+\fr{u}^X.
    \end{equation}

    \item (Kostant). Previously we saw that $\Ad(U^X)H\subset H+\fr{u}^X$. Indeed, we have more: The two 
    sets are equal $\Ad(U^X)H= H+\fr{u}^X$ and for any $W\in\fr{u}^X$ there exist an unique $Z\in\fr{u}^X$ such that 
    $\Ad_{\exp(Z)}(H)=H+W$.\par 
    Let $\fr{u}^X=\bigoplus_{i=1}^m\fr{g}_i^X$. We construct $Z$ inductively as follows. 
    Let $Z_0=0$ and for $j\geq1$ define 
    \begin{equation}
        \begin{split}
            &Z_{j+1}=Z_j+\frac{1}{j+1}p_{j+1}(\Ad_{\exp(Z_j)}(H)-H-W),\\
            &\quad(p_j=\fr{u}^X\to\fr{g}_j^X\;\text{ is the projection}).
        \end{split}
    \end{equation}
    We claim that this sequence satisfies two conditions:
    \begin{equation}
        \begin{split}
            &(1)\; Z_j\in\bigoplus_{i=1}^j\fr{g}_i^X\\
            &(2)\;\Ad_{\exp(Z_j)}(H)-H-W\in\bigoplus_{i=j+1}^m\fr{g}_i^X
        \end{split}
    \end{equation}
    for all $j=1,\dots, m$. 
    If these two conditions are satisfied, then it'z clear that $Z=Z_m$ is the solution of the problem.
    Then we proceed by induction. For $j=1$, we have $Z_1=-W_1\in\fr{g}_1^X$ by definition of the projection. Moreover
    \begin{equation}
        \Ad_{\exp(Z_1)}(H)-H-W\in (\ad(Z_1)H-W)+\sum_{i\geq2}\fr{g}_i^X.
    \end{equation}
    This follows from the fact that $\ad(Z_1)^j:\fr{g}_0^X\to\fr{g}_j^X$. Moreover, by definition
    \begin{equation}
        \ad(Z_1)H-W=-\ad(W_1)H-W=W_1-W\in\bigoplus_{j\geq2}\fr{g}_j^X.
    \end{equation}
    Now, assume the result is true for some $j\geq1$. It's clear from the definition that 
    $Z_{j+1}\in \bigoplus_{1}^{j+1}\fr{g}_i^X+\fr{g}_{j+2}^X=\bigoplus_{1}^{j+2}\fr{g}_i^X$. Now, by 
    Baker-Campbel-Hausdorff formula we have the sequence of equalities modulo $\bigoplus_{i\geq j+2}\fr{g}_i^X$:
    \begin{equation}
        \begin{split}
            &\Ad_{\exp(Z_{j+1})}H-H-W\\
            &=\Ad_{\exp(Z_j)}\Ad_{\exp(Z_{j+1}^\prime)/(j+1)}H-H-W\\
            &=\Ad_{\exp(Z_j)}\left(\Ad_{\exp(Z_{j+1}^\prime)/(j+1)}H-H\right)+\Ad_{\exp(Z_j)}H-H-W
        \end{split}
    \end{equation}

Now, since $\bigoplus_{\geq j+2}\fr{g}_i^X\subset\Ker (p_{j+1})$, we get 
\begin{equation}
    \begin{split}
        &p_{j+1}\Ad_{\exp(Z_j)}\left(Ad_{\exp(Z_{j+1}^\prime)/(j+1)}H-H\right)\\
        &=p_{j+1}\Ad_{\exp(Z_j)}(\frac{[Z_{j+1}^\prime,H]}{j+1})\quad([H,Z_{j+1}^\prime]=(j+1)Z_{j+1}^\prime)\\
        &=-(j+1)\frac{Z_{j+1}^\prime}{j+1}\quad(Z_{j+1}^\prime\in\fr{g}_{j+1}^X)\\
        &=-Z_{j+1}^\prime
    \end{split}
\end{equation}
Therefore, we have
\begin{equation}
    \Ad_{\exp(Z_{j+1})}H-H-W=-Z_{j+1}^\prime+\Ad_{\exp(Z_j)}H-H-W\;\text{ mod }\bigoplus_{\geq j+2}\fr{g}_i^X
\end{equation}
which implies that $p_{j+1}(\Ad_{\exp(Z_{j+1})}H-H-W)=0$. The induction hypothesis yields 
$\Ad_{\exp(Z_j)}H-H-W\in\bigoplus_{i=j+1}^{m}\fr{g}_i^X$. From this, we conclude the induction step.

\item (Kostant). If $\{Y,H,X\},\{Y^\prime,H^\prime,X\}$ are two standard triples with the same nilpositive element. Then 
there exist $x\in U^X$ such that $\Ad_x(Y)=Y^\prime,\Ad_x(H)=H^\prime$ and $\Ad_x(X)=X$.\par
We have $[Y-Y^\prime,X]=H-H^\prime\in[\fr{g},X]$ while $[H-H^\prime, X]=2X-2X=0$. So 
\begin{equation}
    H-H^\prime\in\fr{g}^X\cap[X,\fr{g}]=\fr{u}^X.
\end{equation} 
Therefore, by the previous theorem from Kostant we have $x\in U^X$ such that $\Ad_x(H)=H+(H^\prime-H)=H^\prime$. In particular
$x=\exp(Z),\;Z\in\fr{g}^X$, so $\Ad_x(X)=\exp(\ad(Z))X=X$ and $\{\Ad_x(Y),H^\prime, X\}$ is a standard triple. Since 
$\{Y^\prime,H^\prime,X\}$ is a standard triple with the same neutral and nilpositive element, we must have $\Ad_x(Y)=Y^\prime$.

\item (Mal'cev). If $\{Y,H,X\},\{Y^\prime,H,X^\prime\}$ are two standard triples with the same neutral element. Then, these 
two triples are conjugate.\par 
Let $\fr{g}=\bigoplus_{i\in\mathbb{Z}}\fr{g}_i$ be the decomposition of $\fr{g}$ into $\ad(H)-$eigenspaces.
Also, define $\mathcal{P}=\{Z\in\fr{g}_2:\fr{g}^Z\cap\fr{g}_{-2}\}$ which contains the candidates for nilpositive elements 
for $\fr{sl}_2-$triples containing $H$ as its neutral element. In this case, the group $(G_{\ad}^H)^\circ$ acts in 
$\mathcal{P}$. 
\begin{dang}
    \begin{remark}
        The subgroup $G_{\ad}^H\subset G_{\ad}$ is defined as the group integrating $\fr{g}^H$ and 
        $(G_{\ad}^H)^\circ$ is it connected component containing the identity. The advantage to take 
        the connected component of the identity is that 
        \begin{equation}
            (G_{\ad}^H)^\circ=\exp(\ad(\fr{g}^H)).
        \end{equation}
    \end{remark}
\end{dang}
In fact, we have for $x\in (G_{\ad}^H)^\circ$
\begin{equation}
    \begin{split}
        Z\in\mathcal{P}&\iff \fr{g}^Z\cap\fr{g}_{-2}=0\\
        &\iff x(\fr{g}^Z\cap\fr{g}_{-2})=0 \quad(x \text{ is an isomorphism})\\
        &\iff \fr{g}^{x\cdot Z}\cap x(\fr{g}_{-2})=0\quad(x \text{ is an isomorphism})\\
        &\iff \fr{g}^{x\cdot Z}\cap \fr{g}_{-2}=0\\
        &\quad(\fr{g}^H=\fr{g}_0, \text{ and }\exp(\ad(\fr{g}_0))\fr{g}_{-2}=\fr{g}_{-2})\\
        &\iff x\cdot Z\in\mathcal{P}.
    \end{split}
\end{equation}

The result will follow if we show that $\mathcal{P}$ is connected, while each orbit is clopen. For the first claim, 
we will prove more: $\mathcal{P}$ is path connected. Define the map $T:\fr{g}_{2}\to\Hom(\fr{g}_{-2},\fr{g}_0)$ 
given by $T(Z)=\ad(Z)|_{\fr{g}_{-2}}$. Now we see that 
\begin{equation}
    \begin{split}
        Z\in\mathcal{P}&\iff \fr{g}^Z\cap\fr{g}_{-2}=0\\
        &\iff \text{Ker}T(Z) = 0\\
        &\iff \dim(\text{Im}T(Z))=\dim([Z,\fr{g}_{-2}])=\dim(\fr{g}_{-2})\\
        &\iff T(Z) \text{ has full rank}
    \end{split}
\end{equation}
However, $T(Z)$ is a $\dim(\fr{g}_{-2})\times \dim(\fr{g}_0)-$matrix and the condition of maximal rank is complementary 
to the condition that each of its $\dim(\fr{g}_{-2})\times \dim(\fr{g}_{-2})-$minor is singular. Moreover, since each entry of 
$T(Z)$ is linear in $Z$, this last condition is equivalent to the vanishing of a set of polynomials in $\fr{g}_2$.
So we see that 
\begin{equation}
    \mathcal{P}=\fr{g}_2\setminus 
    \{\det(M(Z))=0: M(Z)\in M_{\dim(\fr{g}_{-2})\times \dim(\fr{g}_{-2})}(\bbC) \text{ a minor of }T(Z)\}.
\end{equation}
So $\mathcal{P}$ is Zariski open and path connected. Let $v,w\in\mathcal{P}$, then for each minor $M(Z)$ of $T(Z)$ there 
exist a finite number of complex numbers 
such that $\det(M(\lambda v+(1-\lambda)w))=0$. Then, there exist a finite collection of complex numbers 
$\{\lambda_1,\dots,\lambda_k\}$ such that $\det(M(\lambda v+(1-\lambda)w))=0$ for each 
$\dim(\fr{g}_{-2})\times \dim(\fr{g}_{-2})-$minor $M(Z)$ of $T(Z)$.
So, we get a continuous path 
$\gamma:[0,1]\to\bbC\setminus\{\lambda_1,\dots,\lambda_k\}$, that composed with $\lambda\mapsto \lambda v+(1-\lambda)w$ yields 
a path connecting $v$ and $w$ in $\mathcal{P}$.\par 

To see that each orbit is open, we will compute dimension of each orbit. We have 
\begin{equation}
    \begin{split}
      \dim( (G_{\ad}^H)^\circ\cdot Z) &= \dim(G_{\ad}^H\cdot Z)\\
      &=\dim(G^H/G_Z^H)\\
      &=\dim(\fr{g}^H)-\dim(\fr{g}^H\cap\fr{g}^Z)\\
      &=\dim(\fr{g}_0)-\dim(\fr{g}^Z\cap\fr{g}_0).
    \end{split}
\end{equation}
Moreover, by the nondegeneracy and invariance of the Killing form, we get 
\begin{equation}
    0=K([\fr{g}^Z\cap\fr{g}_0,Z],\fr{g}_{-2})=K(\fr{g}^Z\cap\fr{g}_0,[Z,\fr{g}_{-2}]).
\end{equation}
So $[Z,\fr{g}_{-2}]\subset (\fr{g}^Z\cap\fr{g}_0)^\perp$ and since $K|_{\fr{g}^H\times\fr{g}^H}$ is nondegenerate, we get 
\begin{equation}
    \begin{split}
        &\dim(\fr{g}^Z\cap\fr{g}_0)^\perp+\dim(\fr{g}^Z\cap\fr{g}_0)=\dim(\fr{g}_0)
    \end{split}
\end{equation}
Therefore
\begin{equation}
    \dim( (G_{\ad}^H)^\circ\cdot Z)\geq \dim(\fr{g}^Z\cap\fr{g}_0)^\perp\geq\dim([Z,\fr{g}_{-2}])=\dim(\fr{g}_{-2})
\end{equation}
and $(G_{\ad}^H)^\circ$ is of maximal dimensional in $\mathcal{P}$, so its open in $\mathcal{P}$. Now, since 
$\mathcal{P}$ is a disjoint union of orbits, we see that each orbit is also closed.\qed
\end{itemize}

\section{Weighted Dynkin diagrams}
There is some difference between mine and Collingwood notation:
        \begin{center}
            \begin{tabular}{|c|c|c|}
                \hline
              &  Mine & Collingwood \\
                \hline 
        Roots & $\Delta$ & $\Phi$\\
        Simple roots & $\Pi$ & $\Delta$\\
        Fundamental domain & $D_\Pi$ & $D_\Delta$\\
        \hline
            \end{tabular}
        \end{center}


De idea here is to classify nilpotent orbits by distinguished semisimple orbits. Here we will fix 
a Cartan subalgebra $\fr{h}\subset \fr{g}$ with associated decomposition $\fr{g}=\bigoplus_{\alpha\in\Delta}\fr{g}_\alpha$.

\begin{equation}
    \mathcal{O}_X\overset{\text{Jacobson-Morozov}}{\longrightarrow}\{Y,H,X\}\to\mathcal{O}_H.
\end{equation}
So we will classify semisimple orbits that cames from neutral elements belonging to a $\fr{sl}_2-$triple. 

The idea is the following:
\begin{itemize}
    \item Let $\{Y,H,X\}$ be a $\fr{sl}_2-$triple with $H$ being the neutral semisimple element. Since $H$ is semisimple, it 
    belongs to some Cartan algebra $\fr{h}_X$ which must be conjugate to $\fr{h}$. So we can take $H\in\fr{h}$ because the 
    adjoint action does not change orbits. Now, we can conjugate $H$ again to take it in the ``fundamental'' 
    domain $D_\Pi$ (why ?).
    \item Let 
    \begin{equation}
        \begin{split}
            &\fr{n}=\bigoplus_{\alpha\in\Delta_+}\fr{g}_\alpha,\\
            &\overline{\fr{n}}=\bigoplus_{\alpha\in\Delta_-}\fr{g}_\alpha.
        \end{split}
    \end{equation}
    Then we have the Borel subalgebra $\fr{b}=\fr{h}\oplus \fr{n}$ and the negative Borel 
    subaglebra $\overline{\fr{b}}=\fr{h}\oplus \overline{\fr{n}}$. Recall the definition of the 
    \textit{fundamental domain}:
    \begin{equation}
        D_\Pi=\{H\in\fr{h}: \text{Re}(\alpha(H))\geq0,\;\text{Im}(\alpha(H))=0,\;\forall\alpha\in\Pi\}.
    \end{equation}
    
    \item We claim that if $H\in D_\Pi$, then it can be completely determined by a map in
    \begin{equation}
        \{0,1,2\}^\Pi
    \end{equation}
    called \textit{weighted Dynkin diagram}. First, 
    remark that $\fr{h}^*=\text{Span}(\Delta)=\text{Span}(\Pi)$. So $\alpha(H)=0\;\forall\alpha\in \Pi$ if and 
    only if $H=0$. Now, we have the $\ad(H)-$eigenspace decomposition:
    \begin{equation}
        \fr{g}=\bigoplus_{i\in\mathbb{Z}}\fr{g}_i.
    \end{equation}
    By the representation theory of $\fr{sl}_2$ we see that $\Ker\;\ad(X)$ is spanned only by highest weight spaces:
    \begin{equation}
        \fr{g}^X=\bigoplus_{i\geq0}\fr{g}_i^X,\;\fr{g}_i^X=\fr{g}^X\cap\fr{g}_i
    \end{equation}
    and in a similar way: $\fr{g}^Y=\bigoplus_{i\leq0}\fr{g}_i^Y$.


\begin{dang}
    Suppose that $\fr{g}$ has the following decomposition as a $\fr{sl}_2-$module:
    \begin{center}
        \begin{tikzcd}
            &&\bullet&&\\
             &\bullet&\bullet&\textcolor{red}{\bullet}_{\text{highest weight}=1}&\\ 
            &\bullet&\bullet&\textcolor{red}{\bullet}_{\text{highest weight}=1}&\\
            \bullet&\bullet&\bullet&\bullet_{\text{weight}=1}&\textcolor{red}{\bullet}_{\text{highest weight}=3}
        \end{tikzcd}
    \end{center}
    Here $\fr{g}_1^X$ has dimension 2, while $\fr{g}_1$ has dimension 3. Moreover $\fr{g}_3=\fr{g}_3^X$.
\end{dang}

    \begin{dang}
        \begin{remark}
            $\fr{g}_i,\; i$ highest weight for the action of $\bbC\langle Y,H,X\rangle$ is just $\fr{g}_i^X$.
        \end{remark}
    \end{dang}

\begin{dang}
    \begin{remark}
        We have $\alpha(H)\in\mathbb{Z}_{\geq0}$, by the representation theory of $\fr{sl}_2$ and the fact that 
        $H\in D_\Pi$.
    \end{remark}
\end{dang}
    
    Now, let $\alpha\in\Pi,\beta \in\Delta_-$, then $\alpha+\beta\in\Delta$ only if  
    if $\alpha+\beta\in\Delta_-$. Otherwise, if $\alpha+\beta\in\Delta_+$, 
    then $\alpha=(\alpha+\beta)-\beta$ will not be simple.\par
    With this observation we see that 
    \begin{equation}
        [X_\alpha, Y]\in \overline{\fr{n}},\;\forall \alpha\in\Pi,
    \end{equation}
    because $Y\in\overline{\fr{n}}$, since $\ad(H)Y=-2Y$ and $H\in D_\Pi$. Then 
    \begin{equation}
        \begin{split}
            &[X_\alpha,Y]=0\Rightarrow X_\alpha\in\fr{g}^Y\Rightarrow \alpha(H)\leq0, \text{ so } \alpha(H)=0\\
            &[X_\alpha,Y]\neq0\Rightarrow \alpha(H)-2\leq0\text{ (by the above observation), so } \alpha(H)\leq 2 
        \end{split}
    \end{equation}
    \item Now, for two nilpotent orbits $\mathcal{O},\mathcal{O}^\prime$ we have associated 
    semisimple elements $H,H^\prime\in D_\Pi$ and the corresponding Dynkin diagrams 
    \begin{equation}
        D(H),D(H^\prime)\in\{0,1,2\}^\Pi.
    \end{equation}
    If $D(H)=D(H^\prime)\neq0$, then $H=H^\prime\neq0$ because of the observation made before: 
    $\alpha(H-H^\prime)=0,\;\forall \alpha\in\Pi\Rightarrow H-H^\prime=0$.
\end{itemize}
\qed

\begin{dang}
    \begin{lemma}
        We have $\#\Pi=\dim(\fr{h}^*)=\text{Rank}(\fr{g})$.
    \end{lemma}
\end{dang}
Proof: We know that $\text{Span}(\Pi)=\text{Span}(\Delta)$, by definition of simple roots. So the result yields provided
we show that the set $\Pi$ is linearly independent. We prove it in two steps:
\begin{itemize}
    \item Claim: If $\alpha,\beta\in\Pi$, then $(\alpha,\beta)\leq0$. In fact, if $(\alpha,\beta)>0$, then 
    \begin{equation}
        p-q=2\frac{(\alpha,\beta)}{(\alpha,\alpha)}>0.
    \end{equation}
    This means that $p>q\geq0$, so $p\geq1$ and $\beta-\alpha\in\Delta\Rightarrow \beta-\alpha\in\Delta_-$ and 
    $\alpha-\beta\in\Delta_+$, which is a contradiction with the fact that $\alpha=(\alpha-\beta)+\beta$ is simple.

    \item Now suppose that $\sum_{\alpha\in\Pi}a_\alpha\alpha=0$. Then if at least one of the sets  
    $I=\{\alpha:a_\alpha>0\}$ and 
    $J=\{\alpha:a_\alpha<0\}$ is nonempty, then we can write this linear dependence as 
    \begin{equation}
       \gamma= \sum_{\alpha\in I}a_\alpha\alpha = \sum_{\alpha\in J}(-a_\alpha)\alpha \in \Delta_+\cup\{0\}.
    \end{equation}
    Therefore 
    \begin{equation}
            (\gamma,\gamma)=\sum_{\alpha\in I}\sum_{\beta\in J}(-a_\alpha a_\beta)(\alpha,\beta)\leq0.
    \end{equation}
    So $(\gamma,\gamma)=0$ and $\gamma=0$. However, from this we deduce that for $\beta\in\Pi$ 
    \begin{equation}
        0=(\gamma,\beta)=\sum_{\alpha\in I}a_\alpha(\alpha,\beta),\; (\alpha,\beta)\leq0, a_\alpha\geq0.
    \end{equation}
    This forces $(\alpha,\beta)=0,\;\forall \alpha\in I$. In the same way we see that 
    $(\alpha,\beta)=0,\;\forall \alpha\in J$, so 
    \begin{equation}
              (\alpha,\beta)=0,\;\forall\alpha\in I\cup J.   
    \end{equation}
    So $\alpha\in I\cup J \neq \varnothing\Rightarrow (\alpha,\alpha)=0$ and $\alpha=0$, which is a contradiciton.\qed
\end{itemize}


\begin{dang}
    \begin{lemma}
        There is an isomorphism: 
\begin{equation}
     W \overset{\cong}{\longrightarrow} N_{G_\ad}(\fr{h})/C_{G_\ad}(\fr{h}).
\end{equation}
    \end{lemma}
\end{dang}
Proof: We will construct this isomorphism in several steps:
\begin{itemize}
    \item We can transpose the action of $W$ to $\fr{h}$ vi the contragradient representation:
    \begin{equation}
        \langle w\cdot h, \alpha\rangle = \langle h,w^{-1}(\alpha)\rangle.
    \end{equation}
    Write $s_\alpha$ for the action of $r_\alpha$ transposed to $\fr{h}\cong\fr{h}^{**}$. Then we have 
    \begin{equation}
        \begin{split}
           \langle s_\alpha(H),\lambda \rangle&= \langle H,r_\alpha(\lambda)\rangle\\
            &=\langle H, \lambda -2\frac{(\alpha,\lambda)}{(\alpha,\alpha)}\alpha\rangle\\
            &=\langle H, \lambda-2\frac{\lambda(\nu^{-1}(\alpha))}{(\alpha,\alpha)}\alpha\rangle\\
            &=\langle H,\lambda\rangle -2\frac{1}{(\alpha,\alpha)}\langle H,\alpha\rangle
            \langle \nu^{-1}(\alpha),\lambda \rangle\\
            &=\langle H-2\frac{\langle H,\alpha\rangle}{(\alpha,\alpha)}\nu^{-1}(\alpha),\lambda\rangle\\
            &=\langle H-\langle H,\alpha\rangle H_\alpha,\lambda\rangle.
        \end{split}
    \end{equation}
    So $s_\alpha(H)=H-\alpha(H)H_\alpha$, defines the action of $W\cong\langle s_\alpha:\alpha\in\Delta\rangle$ on
    $\fr{h}$.
    \item Define $\theta_\alpha = \exp(\ad(E_\alpha))\exp(\ad(-F_\alpha))\exp(\ad(E_\alpha))\in\Aut(\fr{g})$, where 
    $\{F_\alpha,H_\alpha, E_\alpha\}$ is the standard $\fr{sl}_2-$triple, i.e., $(E,F)\in\fr{g}_\alpha\times\fr{g}_{-\alpha}$ is 
    such that $K(E,F)=2/K(\alpha,\alpha)$ and $H_\alpha = 2\nu^{-1}(\alpha)/K(\alpha,\alpha)$. 
    We can show that $\theta_\alpha|_{\fr{h}}=s_\alpha$. So each 
    reflection $s_\alpha$ can be extended to an automorphism of $\fr{g}$.
    \begin{dang}
        \begin{remark}
            If $\delta\in\text{Der}(\fr{g})$ is a nilpotent derivation, then $\exp(\delta)\in\Aut(\fr{g})$.
        \end{remark}
    \end{dang}
    \item It's a fact that $\theta_\alpha \in \Ad(G)\subset\Aut(\fr{g})$. This follows by the identity 
    $\Ad\circ\exp(Z)=\exp\circ\ad(Z)$. Also, we have  
    $\langle \theta_\alpha:\alpha\in\Delta \rangle=N_{G_{ad}}(\fr{h})$. So we have a surjective map 
    \begin{equation}
        W\to N_{G_{ad}}(\fr{h}): r_\alpha\mapsto \theta_\alpha
    \end{equation}
    that induces the desired isomorphism.

    \begin{center}
        \textcolor{red}{We need to be more careful here. See section~(\ref{section: isom of Weyl groups})}
    \end{center}
\end{itemize}

Now lets do an example in $\fr{sl}_n$. In this case, the root space is 
\begin{equation}
    \Delta =\{L_i-L_j: i\neq j\}\subset V=\text{Span}(L_i)_{1\leq i\leq n}=\mathbb{R}^n.
\end{equation}
with respect to the decomposition
\begin{equation}
    \fr{sl}_n=\bigoplus_{i>j}\bbC E_{ij}\oplus\fr{h}\oplus \bigoplus_{i<j}\bbC E_{ij}
\end{equation}
where $\fr{h}=\fr{sl}_n\cap\{\text{diagonal matrices}\}$. Of course $\fr{g}_{L_i-L_j}=\bbC E_{ij}$ as can be seen by the
equation $\ad(H)E_{ij}=(L_i-L_j)(H)E_{ij}$.\par
We define a linear functional $f: V\to\mathbb{R}: L_i\mapsto -i$. So, the positive roots with respect to this choice 
of functional is $\Delta_+=\{L_i-L_j:i<j\}$. The Cartan triples are $\{E_{ji},H_{i,j},E_{ij}\}$ where 
$H_{i,j}=E_{ii}-E_{jj}$. So, the Cartan subalgebra $\fr{h}$ lies in the borel set of upper triangular matrices in 
$\fr{sl}_n$. Said this, the set of simple roots is 
\begin{equation}
    \Pi =\{L_1-L_2,L_2-L_3,\dots, L_{n-1}-L_n\}.
\end{equation}
In fact, all other positive roots can be writen as a linear combination of these roots. The fact that this spanning set 
has $n-1=\text{Rank}(\fr{sl}_n)$ elements, shows that they must be the simple roots. Now, 
$H=\text{Diag}(\lambda_1,\dots,\lambda_n)\in D_\Pi$ only if $\alpha(H)\geq0\;\forall\alpha\in\Pi$, i.e,
\begin{equation}
    \lambda_1\geq\lambda_2\geq\cdots\geq\lambda_n.
\end{equation}
This define the $\Pi-$dominant set in $\fr{h}$. Remember that $[e,f]=K(e,f)\nu^{-1}(\alpha)$ for 
$(e,f)\in\fr{g}_{-\alpha}\times\fr{g}_\alpha$. So, if $h_\alpha=K(e,f)\nu^{-1}(\alpha)$ we get 
\begin{equation}
    [h_\alpha,e]=K(e,f)(\alpha,\alpha)e=2e\iff \frac{2}{(\alpha,\alpha)}=K(e,f).
\end{equation}
So, the fact that $\{E_{ji},H_{i,j},E_{ij}\}$ forms a $\fr{sl}_2-$triple shows that $H_{i,j}=H_{L_i-L_j}$ in the sence 
of the definition of a coroot $h_\alpha=\frac{2}{(\alpha,\alpha)}\nu^{-1}(\alpha)$. Now, we can use the known expression
for the action of the Weyl group on coroots:
\begin{equation}
    \begin{split}
        s_{L_i-L_j}(H)&=H-(\lambda_i-\lambda_j)H_{L_i-L_j}\\
        &=H-(\lambda_i-\lambda_j)H_{i,j}\\
        &=\text{Diag}(\lambda_1,\dots,\lambda_k-\lambda_k+\lambda_l,\dots,\lambda_l-\lambda_l+\lambda_k,\dots\lambda_n)\\
        &=\text{Diag}(\lambda_1,\dots,\hat{\lambda_l},\dots,\hat{\hat{\lambda_k}},\dots,\lambda_n),\\ 
        &\quad\quad\hat{}=k\text{th position}, \hat{\hat{}}=l\text{th position} \;(k<l).
    \end{split}
\end{equation}
So the transposed action of the weyl group on $\fr{h}$ permutes elements of the diagonal. We know that this action comes 
from elements of $G_{\ad}$, so we are not changing the orbit of $H$ by applying these maps. So, reordering the diagonal elements 
of $H$, we take its conjugate in the domain $D_\Pi$ as a diagonal matrix with decreasing elements on the diagonal.

For $H=\text{Diag}(\lambda_1,\dots,\lambda_n)$, we have the weighted Dynkin diagram:
\begin{equation}
    L_i-L_{i+1} \in\Pi\mapsto \lambda_i-\lambda_{i+1}\in\{0,1,2\}.
\end{equation}

For example, let $n=3$. We have the orbit $\mathcal{O}_{[2,1]}$ with representative $X_{[2,1]}$. Remember that 
\begin{equation}
    X_{[2,1]}=\rho_1(X)\oplus\rho_0(X)=\begin{pmatrix}0&1\\0&0\end{pmatrix}\oplus(0)=E_{1,2}\in\fr{sl}_3.
\end{equation}
The corresponding semisimple element constructed explictly is 
\begin{equation}
    H_{[2,1]}=\rho_1(H)\oplus\rho_0(H)=\begin{pmatrix}1&0\\0&-1\end{pmatrix}\oplus(0)=H_{1,2}.
\end{equation}
Reordering the elements of the diagonal in a decreasing order we get 
$\text{Diag}(1,0,-1)$. Then the associated weighted Dyinkin diagram is
\begin{align*}
    \left\{\begin{aligned}
        &L_1-L_2\mapsto 1,\\
        &L_2-L_3\mapsto 1.
    \end{aligned}\right.
\end{align*}


\section{Regular, subregular and minimal orbits}


\section{Standard slices}
\begin{dang}
    \begin{lemma}\label{lemma:transversality lemma}
        Let $f: M\to \bbR^k,\;g:N\to\bbR^k$ be an imersion and an embedding respectively. Suppose that 
        $x=f(p)=g(q)$ and 
        \begin{equation}
            \bbR^k = f_{*,p}T_pM\oplus g_{*,q}T_qN.
        \end{equation}
        Then, there exist a neighborhoods $U\ni p,V\ni q$ such that $f(U)\cap g(V)=\{x\}$.
    \end{lemma}
\end{dang}
Proof: Remember the smooth strucutre of $M\times N$: If $(\varphi,U)$ and $(\psi,V)$ are charts around $p$ and $q$ respectively. 
Then $(\pi_M^*\varphi,\pi_N^*\psi), U\times V$ is a chart around $(p,q)$. 
If $\varphi=(x_1,\dots,x_m)$ and $\psi=(y_1,\dots, y_n)$, then we will write $\bar{x}_j=\pi_M^*(x_j)$ and 
$\bar{y}_j=\pi_N^*(y_j)$. So $(\bar{x}_1,\dots,\bar{x}_m,\bar{y}_1,\dots,\bar{y}_n)$ is a coordinate chart around $(p,q)$. Under this notation 
it's clear that 
\begin{equation}
    (\pi_M)_{*}(\frac{\partial}{\partial\bar{x}_j})=\frac{\partial}{\partial x_j},\;
    (\pi_N)_{*}(\frac{\partial}{\partial\bar{y}_j})=\frac{\partial}{\partial y_j}
\end{equation}
and that the map 
\begin{equation}
    \left\{\begin{aligned}
        &\frac{\partial}{\partial\bar{x}_j}\mapsto (\frac{\partial}{\partial x_j},0)\\
        &\frac{\partial}{\partial\bar{y}_j}\mapsto (0,\frac{\partial}{\partial y_j})
    \end{aligned}\right.
\end{equation}
defines an isomorphism $T_{(p,q)}(M\times N)\to T_pM\times T_qN$.\par 
Now, let go to the proof. Define the map $\phi:M\times N\to \bbR^k:(p,q)\mapsto f(p)-g(q)$, that is 
$\phi=f\circ\pi_M-g\circ\pi_N$. We have $\phi_*(\frac{\partial}{\partial\bar{x}_j})=f_*(\frac{\partial}{\partial x_j})$ and 
$\phi_*(\tfrac{\partial}{\partial\bar{y}_j})=-g_*(\tfrac{\partial}{\partial y_j})$. So 
\begin{equation}
    \begin{split}
        &\phi_{*,(p,q)}(\sum_ia_i\frac{\partial}{\partial\bar{y}_i}+\sum_jb_j\frac{\partial}{\partial y_j})=0\iff\\
        &f_{*,p}(\sum_ia_i\frac{\partial}{\partial x_i})-g_{*,q}(\sum b_j\frac{\partial}{\partial y_j})=0\iff \\
        &f_{*,p}(\sum_ia_i\frac{\partial}{\partial x_i})=g_{*,q}(\sum_jb_j\frac{\partial}{\partial y_j})=0\,\;(\text{ by transversality})\iff\\
        &a_i=0,\;\forall i\;\text{ and }b_j=0,\;\forall j\quad(f_*,g_*\text{ are injective})
    \end{split}
\end{equation}
This proves that $\phi_*$ is injective and by dimensionality it must be an isomorphism as well. Therefore $\phi$ is a 
local diffeomorphism around $(p,q)$. So there exist $U\times V\ni (p,q)$ such that $\phi|_{U\times V}$ is a diffeomorphism onto its image. 
In particular $(\phi|_{U\times V})^{-1}(0)=(p,q)$ and $f(U)\cap g(V)=\{x\}$.

\begin{dang}
    \begin{theorem}
        Let $\fr{g}$ be a semisimple Lie algebra, $X\in\fr{g}$ a nilpotent element and 
        $\fr{a}=\bbC\langle Y,H,X\rangle$ a $\fr{sl}_2-$triple containing $X$ as its nilpositive element. Then we have 
        \begin{equation}
            \begin{aligned}
             &(1)\quad  T_X\mathcal{O}_X \oplus T_X(X+\fr{g}^Y)= [\fr{g},X]\oplus\fr{g}^Y = \fr{g}\\
             &(2)\quad  \mathcal{O}_X\cap(X+\fr{g}^Y)=\{X\}
            \end{aligned}          
        \end{equation}
    \end{theorem}
\end{dang}
Proof: It's a fact that the inclusion $\mathcal{O}_X\hookrightarrow\fr{g}$ is an imersion and 
$T_X\mathcal{O}_X=\{\ad(X)Z:Z\in\fr{g}\}=[\fr{g},X]$ with the identification $T_X\mathcal{O}_X\cong i_*(T_X\mathcal{O}_X)$. The 
equation $[\fr{g},X]\oplus\fr{g}^Y = \fr{g}$ is a consequence of the representation theory of $\fr{sl}_2$. Now, by exponentiating the relation  
$\ad(zH)X=2zX$ (viewed as elements of $\fr{sl}_2$) we get 
\begin{equation}
   \exp(\ad(zH))X=(\exp z)^2X,\;z\in\bbC.
\end{equation}
which is equivalent to 
\begin{equation}
    \Ad_{\left(\begin{smallmatrix}t&0\\0&\;\;t^{-1}\end{smallmatrix}\right)}X=t^2X,\;t\in\bbC.
\end{equation}
The last equation, can be seen as an equation in $\fr{g}$. To do this integrate the Lie algebra morphism $\fr{sl}_2\to\fr{a}\subset \fr{g}$ to 
a Lie group homomorphism $\gamma:SL_2(\bbC)\to G$. Finally, restricts $\gamma$ to the maximal torus 
$T=\{\left(\begin{smallmatrix}t&0\\0&\;\;t^{-1}\end{smallmatrix}\right):t\in\bbC\}$. Since $\gamma_*$ is a Lie algebra morphism, we 
can do all the calculations made before in $\fr{g}$. The last equation is writen as $\Ad_{\gamma(t)}X=t^2X$ provided we identify 
$T\cong \bbC^*$, that is we write $\gamma(t)$ for $\gamma\left(\begin{smallmatrix}t&0\\0&\;\;t^{-1}\end{smallmatrix}\right)$.\par 
We observe that $t^2\Ad_{\gamma(t^{-1})}X=X$ by construction. Moreover, if $zX\in\mathcal{O}_X$ for any $z\in X$ because 
$\exp(z\ad(H))X=\exp(z)^2X$. So $t^2\Ad_{\gamma(t^{-1})}$ preserves $\mathcal{O}_X$ as well. Furthermore if $Z\in\fr{g}^Y$, then 
\begin{equation}
    [Y,t^2\Ad_{\gamma(t^{-1})}Z]=\Ad_{\gamma(t^{-1})}[t^2\Ad_{\gamma(t)}Y,Z]=\Ad_{\gamma(t^{-1})}[Y,Z]=0.
\end{equation} 
The last equation can be seen by exponentiating the relation $\ad(H)Y=-2Y$. In conclusion, we see that the one parameter group 
$\bbC^*$ preserves $\mathcal{O}_X\cap(X+\fr{g}^Y)$ through the action $t^2\Ad_{\gamma(t^{-1})}$. Now, observe that 
$\fr{g}^Y\subset \bigoplus_{i\leq0}\fr{g}_i,\;\fr{g}_i=E_i^{\ad(H)}$. For $Z\in\fr{g}_i$ we have 
$t^2\Ad_{\gamma(t^{-1})}Z=t^{2-i}Z$. These equations shows that we can send 
\begin{equation}
    \lim_{t\to 0}t^2\Ad_{\gamma(t^{-1})}(X+Z)=X+\lim_{t\to0}t^2\Ad_{\gamma(t^{-1})}Z= X.
\end{equation}

So, if there exist $Z\in\fr{g}^Y\setminus0$ such that $X+Z\in\mathcal{O}_X$, 
this will contradict the lemma~(\ref{lemma:transversality lemma}), because the family 
$\{t^2\Ad_{\gamma(t^{-1})}(X+Z)\}_{t\to0}$ will be arbitrarily close, but different to $X$.\qed 
    \begin{example}
        Let $\fr{g}=\fr{sl}_3$ with $X=E_{12}$. We have the $\fr{sl}_2-$triple $\{E_{21},H_{1,2},E_{12}\}$ and 
        $\fr{sl}_3$ decomposes as $V_{(0)}\oplus V_{(1)}^{\oplus2}\oplus V_{(2)}$ where 
        \begin{equation}
            \begin{aligned}
                &V_{(0)}=\text{Ker}(L_1-L_2)=\bbC\text{Diag}(1,1,-2)\\
                &V_{(1)}=\bbC E_{23}\oplus\bbC E_{13}\\
                &V_{(1)}=\bbC E_{31}\oplus\bbC E_{32}\\
                &V_{(2)}=\bbC E_{21}\oplus \bbC H_{(1,2)}\oplus\bbC E_{12}
            \end{aligned}
        \end{equation}
Now we observe that 
\begin{equation}
    X+\fr{g}^Y=\left\{
        \begin{pmatrix}
            a&1&0\\b&a&c\\d&0&-2a
        \end{pmatrix}:
        a,b,c,d\in\bbC
    \right\}.
\end{equation}
Now, let $W\in X+\fr{g}^Y$, then $W\in\mathcal{O}_X$ only if it's conjugate to $X=E_{12}$. So $W$ must be a rank 1 one matrix. 
This means that all order 2 minors must vanish. Therefore
\begin{equation}
    \begin{aligned}
        &\det\begin{pmatrix}a&1\\b&a\end{pmatrix}=a^2-b=0\\
        &\det\begin{pmatrix}\;\;a&0\\\;\;\;0&-2a\end{pmatrix}=-2a^3=0\\
        &\det\begin{pmatrix}1&0\\a&c\end{pmatrix}=c=0\\
        &\det\begin{pmatrix}a&1\\d&0\end{pmatrix}=d=0
    \end{aligned}
\end{equation}
So $a=b=c=d=0$ and $W=E_{12}=X$.
\end{example}

\section{An isomorphism of Weyl groups}\label{section: isom of Weyl groups}
Notation for this section:
\begin{itemize}
    \item $\fr{g}$ a complex semisimple Lie algebra and $\fr{g}_0$ its real form, i.e., $\fr{g}=\fr{g}_0\otimes\bbC$.
    \item $\fr{h}\subset \fr{g}$ a Cartan subalgebra.
    \item $G$ is the compact Lie group associated to the real form $\fr{g}_0$.
    \item $T\subset G$ a maximal torus and $X^*(T)=\{T\overset{C^\infty}{\to}\bbC^*\}$ the set of its complex characters. 
    \item $N(T)$ is the normalizer of $T$ in $G$.
    \item We will write $W_T=N(T)/T$ and $W$ for the Weyl group associated to the root system of $\fr{g}$.
\end{itemize}
We have the following monomorphism: 
    \begin{equation}\label{equation: embedding of X(T) into a vector space}
        X^*(T)\hookrightarrow V_{\bbR}=\Hom_{\bbR}(\text{Lie}(T),i\bbR):\alpha\mapsto d\alpha,
    \end{equation}
where we identify $\text{Lie}(\mathbb{S}^1)=i\bbR$.
    
\begin{remark}
    A maximal torus $T$, in a Lie group, is a maximal (with respect to inclusion) compact, connected, abelian subgroup 
    of $G$. We can show that in this circunstance we have identifications:
    \begin{equation}
        \begin{aligned}
            &T=(\bbR/\bbZ)^r\cong\underset{r\text{ times}}{\underbrace{\mathbb{S}^1\times\cdots\times\mathbb{S}^1}},\\
            &\bbZ^r\overset{\cong}{\to} X^*(T):(k_1,\dots,k_r)
        \mapsto [(e^{i\theta_1},\dots,e^{i\theta_r})\mapsto \prod_j e^{i\theta_jk_j}].
        \end{aligned}
    \end{equation}
    Later we will see that $r=\text{Rank}(\fr{g})$.
\end{remark}
By the previous equation~(\ref{equation: embedding of X(T) into a vector space}), we get an isomorphism
\begin{equation}
    \psi:\bbR\otimes_{\bbZ}X^*(T)\overset{\cong_{\bbR}}{\longrightarrow} 
    V_{\bbR}:\alpha\mapsto d\alpha.
\end{equation}


Now we build the diagram, where the dashed line represents the passage to the ``dual''.
 \begin{center}
    \begin{tikzcd}
        T\arrow[d, dashed]&\text{Lie}(T)\arrow[l, two heads, swap,"\exp"]\arrow[d, dashed]\\
        \bbR\otimes X^*(T)\arrow[r,"\psi"]&V_{\bbR}
    \end{tikzcd}
\end{center}
There is a natural action of $W_T$ on this diagram. Let $w\in W_T$, then we have:
\begin{itemize}
    \item Conjugation $c_w:T\to T:t\mapsto wtw^{-1}$
    \item The associated dual $c_w^*$ on $X^*(T)$, which takes a character $\chi$ and gives 
    $c_w^*(\chi)=\chi\circ c_{w^{-1}}$
    \item The adjoint action $\Ad_w$ on $\text{Lie}(T)$
    \item The co-adjoint action $\Ad_w^*$ on $V_{\bbR}$
\end{itemize}

We know that $c_w(\exp(H))=\exp(\Ad_w(H))$. In addition, we verify 
\begin{equation}
    d(c_w^*\alpha)=d(\alpha\circ c_w^{-1})=d\alpha\circ\Ad_{w^{-1}}=\Ad_w^*(d\alpha)
\end{equation}
which shows that $\psi$ intertwines the actions of $W_T$ on $V_{\bbR}$ and $\bbR\otimes X^*(T)$. Therefore, 
the maps on the previous diagram intertwines all the actions of $W_T$. Now we observe that 
$\text{Lie}(T)\subset \fr{g}_0$ is maximal abelian subalgebra. Moreover, in the complexficiation, we get that 
$\fr{h}=\text{Lie}(T)_{\bbC}\subset\fr{g}$ is a maximal toral subalgebra (MTS), so it's a Cartan subalgebra.\par 

Lets $\fr{p}$ be a complementary subspace to $\text{Lie}(T)$ in $\fr{g}_0$. It's a fact that 
$\Ad_T$ restricts to $\fr{p}$, so the decomposition $\fr{g}=\fr{h}\oplus \fr{p}_{\bbC}$ is fixed 
by the complexfied action of $\Ad_T$ on $\fr{g}$. Let 
\begin{equation}
    \fr{p}_{\bbC}=\bigoplus_{\alpha\in \Phi}V_\alpha,\; \Phi=\{\alpha\in X^*(T):V_\alpha\neq0\} 
\end{equation}
be the corresponding decomposition into $T-$irreducible modules for its adjoint action. Now, consider the complexified 
map $\psi_{\bbC}:\bbC\otimes_{\bbZ}X^*(T)\to \Hom_{\bbC}(\fr{h},\bbC)=\fr{h}^*$. We claim that 
$\psi_{\bbC}(\Phi)=\Delta$, where $\Delta$ is the set of roots induced by the Cartan subalgebra $\fr{h}$. 
In fact, if $X_\alpha\in V_\alpha$ and $H\in\text{Lie}(T)$, then
\begin{equation}
    \Ad_{\exp(tH)}X_\alpha=\exp(\ad(tH))X_\alpha\overset{\tfrac{d}{dt}\big|_0}{\Longrightarrow} 
    \ad(H)X_\alpha=d\alpha(H)X_\alpha
\end{equation}
and $V_\alpha\subset \fr{g}_{d\alpha}$. The right hand side of the last equation is true for $H\in\fr{h}$ by linear 
extension. The converse inclusion we get by exponentiation of the adjoint action of 
$\fr{h}$ on $\fr{g}$. It follows that $V_\alpha=\fr{g}_{d\alpha}$ 
and $\fr{p}_{\bbC}=\bigoplus_{\alpha\in\Phi}\fr{g}_{d\alpha}$. For simplicity of notation, we will write 
$\fr{g}_\alpha$, instead of $\fr{g}_{d\alpha}$.\par 

Remember that if $f:\Delta\to\bbR^*$ is a linear functional, then we have an induced set of positive roots $\Delta_+$. 
Now if $\alpha\in \Delta_+,\; X_{\alpha}\in\fr{g}_{\alpha},\; Y_{\alpha}\in\fr{g}_{-\alpha}$ and 
$K(X_\alpha,Y_\alpha)=2/(\alpha,\alpha)$, then 
\begin{equation}
    H_\alpha = [X_\alpha,Y_\alpha]= \frac{2}{(\alpha,\alpha)}\nu^{-1}(\alpha)
\end{equation} 
and we get a $\fr{sl}_2-$triple $\fr{sl}_\alpha=\bbC\langle Y_\alpha,H_\alpha,X_\alpha\rangle$. This defines morphisms 
$\fr{sl}_2\bbC\hookrightarrow\fr{g}$ and $di_\alpha:\fr{su}(2)\hookrightarrow\fr{g}_0$, of real forms. Since 
$SU(2)$ is simple connected, we have an associated map of Lie groups $i_\alpha:SU(2)\to G$.

\begin{remark}
    \begin{equation}
        \begin{aligned}
            SU(2)&=\{X\in GL_2(\bbC): A\overline{A}^T=I,\;\det(A)=1\}\\
                 &=\left\{\begin{pmatrix}a&b\\-\overline{b}&\overline{a}\end{pmatrix}: |a|^2+|b|^2=1\right\}\\
       \fr{su}(2)&=\{X\in M_2(\bbC):X+\overline{X}^T=0,\;\Tr(X)=0\}\\
                 &=\bbR\begin{pmatrix}
                 &-1\\1& 
            \end{pmatrix}\oplus
            \bbR\begin{pmatrix}
               i&\\&-i 
            \end{pmatrix}\oplus 
            \bbR\begin{pmatrix}
               &i\\i& 
            \end{pmatrix}
        \end{aligned}
    \end{equation}
\end{remark}

\begin{dang}
    \begin{notation}
        Write $T_\alpha = \Ker(\alpha)\subset T$ and $\fr{su}_\alpha = di_\alpha(\fr{su}(2))$.
    \end{notation}
\end{dang}

\begin{dang}
    \begin{lemma}
        We have
        \begin{itemize}
            \item[1.] $\Ad_w(H)=H,\;\forall H\in\text{Lie}(T_\alpha)=\Ker(d\alpha)$ and $w\in i_\alpha(SU(2))$
            \item[2.] $\Ad_{w_\alpha}(H_\alpha)=-H_\alpha$
            \item[3.] $w_\alpha=i_\alpha(\begin{smallmatrix}&-1\\1&\end{smallmatrix})\in N(T)$
        \end{itemize}
    \end{lemma}
\end{dang}
Proof: Consider the exponential map $\fr{su}(2)\twoheadrightarrow SU(2)$, then if $w\in i_\alpha(SU(2))$, we have
$w=\exp(X)$ for some $X\in \fr{su}_\alpha$. Therefore $\Ad_w(Y)=\exp(\ad(X))Y=Y$ provided we show that $[X,Y]=0$. 
But this is the case, since $[\fr{su}_\alpha,\Ker(d\alpha)]=0$. This concludes the first assertion.\par 

For the second assertion, we observe that 
$A_\alpha=X_\alpha-Y_\alpha=(\begin{smallmatrix}&1\\-1&\end{smallmatrix})$, and 
$\Ad_{w_\alpha} = \exp(\tfrac{\pi}{2}\ad(A_\alpha))$. Define $S_\alpha = X_\alpha+Y_\alpha$, then 
$\ad(X_\alpha)$ acts on $U=\bbR\langle H_\alpha, S_\alpha\rangle$ as the matrix 
$(\begin{smallmatrix}&2\\-2&\end{smallmatrix})$. Therefore
\begin{equation}
  \Ad_{w_\alpha}=\exp(\frac{\pi}{2}\ad(A_\alpha))\big|_U = -\text{Id}.
\end{equation}

Finally, we observe that $\text{Lie}(T)=\bbR iH_\alpha\oplus \Ker(d\alpha)$ because $d\alpha(H_\alpha)\neq0$. Therefore 
\begin{equation}
    T=\langle \underset{S}{\underbrace{\{e^{s iH_\alpha}\}_{s\in\bbR}}}, T_\alpha \rangle.
\end{equation}
However
\begin{equation}
    \begin{split}
        & \left\{\begin{aligned}
        &1\text{th assertion }\Rightarrow c_w(e^H)=e^H,\;\forall H\in\text{Lie}(T_\alpha)\\
        &2\text{th assetion }\Rightarrow c_{w_\alpha}(e^{siH_\alpha})=e^{-siH_\alpha}
    \end{aligned}\right. \\
     &\quad\Longrightarrow\left\{\begin{aligned}
        & w\in C_G(T_\alpha),\;\forall w\in i_\alpha(SU(2))\\
        & w_\alpha\in N(S)
     \end{aligned}\right.
\end{split}
\end{equation}
This proves that $w_\alpha$ normalizes $T$, that is $w_\alpha\in N(T)$.\qed

For the first assertion it follows that if $(d\alpha,d\beta)=0$, then $\Ad_{w_\alpha}^*d\beta=d\beta$. In fact, if
$(d\alpha,d\beta)=0$, then $d\alpha(H_\beta)=0$ and $H_\beta\in\text{Lie}(T_\alpha)$. Therefore 
\begin{equation}
    \Ad_{w_\alpha}(H_\beta)=H_\beta\Rightarrow \Ad_{w_\alpha}^*d\beta=d\beta
\end{equation}
follows by the definition of a coroot and the fact that the isomorphism $\nu:\fr{h}\to\fr{h}^*$, induced by the 
Killing form, intertwines the adjoint and coadjoint actions of $G$.\par 

For the second assertion of the previous lemma we see that $\Ad_{w_\alpha}^*d\alpha=-d\alpha$. In fact, by applying 
the Killing form we see that
\begin{equation}
    \begin{split}
        \Ad_{w_\alpha}(H_\alpha)=-H_\alpha\Rightarrow d\alpha(\Ad_{w_\alpha^{-1}}H)=-d\alpha(H),\;\forall H\in\fr{h}
    \end{split}
\end{equation}
and the assertion follows.\par 

In conclusion, the two previous observations shows that $\Ad_{w_\alpha}^*=r_\alpha$. In fact, for any $\lambda\in V_{\bbR}$, 
we have $\lambda-\tfrac{(\lambda,\alpha)}{(\alpha,\alpha)}\alpha\in \bbR\alpha^\perp$. Therefore 
\begin{equation}
    \Ad_{w_\alpha}^*(\lambda-\tfrac{(\lambda,\alpha)}{(\alpha,\alpha)}\alpha)
    =\lambda-\tfrac{(\lambda,\alpha)}{(\alpha,\alpha)}\alpha,
\end{equation}
which proves that $\Ad_{w_\alpha}^*(\lambda)=\lambda-2\tfrac{(\lambda,\alpha)}{(\alpha,\alpha)}\alpha$. Now let 
$W^\prime=\langle w_\alpha:\alpha\in\Phi \rangle\subset W_T$, then we see that we get an isomorphism
\begin{equation}
    W^\prime \longrightarrow W: w_\alpha\mapsto \Ad_{w_\alpha}^*.
\end{equation}
The next lemma will finnish the proof of the isomorphism.

\begin{dang}
    \begin{lemma}
        We have $W^\prime = W_T$.
    \end{lemma}
\end{dang}
Proof: First we observe that
\begin{equation}
    X_{c_w^*\alpha}=\Ad_w(X_\alpha)\neq0,\;\text{if } \alpha\in\Phi,w\in W_T.
\end{equation}
So $W_T$ acts on $\Phi$. In particular, this proves that $W_T$ is finite. So there exist $g_i\in N(T)\;(i=1,\dots,k)$ such that
$N(T)=\coprod_{i=1}^kg_iT$ and we see that the left cosets $g_iT$ are the connected components of $N(T)$. 
Moreover, since $T$ is the connected component that contains the unity, we see that $T=N(T)^\circ$. Now define 
the Weyl vector $\rho=\tfrac12\sum_{\alpha\in\Delta_+}\alpha$. It's a fact that it belongs to the positive Weyl chamber 
\begin{equation}
    \mathcal{C}_+^\circ = \{\lambda\in V_{\bbR}: (\lambda,\alpha)>0,\;\forall\alpha\in\Delta_+\}.
\end{equation}
Therefore $\rho$ is not fixed by $\Ad_{w_\alpha}^*$ by any $\alpha\in\Phi$. This shows that 
\begin{equation}
   H_0= \nu^{-1}(\rho)=\frac12\sum_{\alpha\in\Delta_+}\nu^{-1}(\alpha)\in\fr{h}.
\end{equation}
is not fixed by any $\Ad_{w_\alpha},\;\alpha\in\Delta$.\par 

Now let $\tilde{w}\in W_T$, we know that $\Ad_{\tilde{w}}(\Phi)=\Phi$, so there exist $w\in W^\prime$ such that 
$\Ad_{w^{-1}}^*\Ad_{\tilde{w}}^*=\Ad^*_{w^{-1}\tilde{w}}$ permutes the elements of $\Delta_+$. Write $n=w^{-1}\tilde{w}$, then 
we see by the previous equation that $\Ad_n$ fixes $H_0$. Therefore 
\begin{equation}
    \begin{aligned}
        &H_0 \subset \text{Lie}(T)\setminus \bigcup_{\alpha\in\Delta}\text{Lie}(T_\alpha)\\
        &\quad\Longrightarrow S = \{e^{tH_0}\}_{t\in\bbR}\subset T \text{ contains a regular element }
        e^{H_0}\notin\bigcup_{\alpha\in\Delta}T_\alpha
    \end{aligned}
\end{equation}

Then we see that $C_G(S)\subset N(T)\overset{C_G(S)\text{ connected}}{\Longrightarrow}C_G(S)\subset N(T)^0=T$. However
$c_n(e^{H_0})=\exp(\Ad_n(H_0))=e^{H_0}$ and $n\in C_G(S)\Rightarrow n\in T$. Therefore $\tilde{w}=w\in W_T$, i.e., 
$W_T=W^\prime$.\qed 

